\documentclass[letter, 10pt]{article}
\usepackage[latin1]{inputenc}
\usepackage[spanish]{babel}
\usepackage{amsfonts}
\usepackage{amsmath}
\usepackage[dvips]{graphicx}
\usepackage{url}
\usepackage[top=3cm,bottom=3cm,left=3.5cm,right=3.5cm,footskip=1.5cm,headheight=1.5cm,headsep=.5cm,textheight=3cm]{geometry}



\begin{document}
\title{Inteligencia Artificial \\ \begin{Large}Primer avance: Problema Generaci�n Procedural de Plataformas\end{Large}}
\author{Bruno Morici}
\date{\today}
\maketitle


%--------------------No borrar esta secci\'on--------------------------------%
%\section*{Evaluaci\'on}

%\begin{tabular}{ll}
%Resumen (5\%): & \underline{\hspace{2cm}} \\
%Introducci\'on (5\%):  & \underline{\hspace{2cm}} \\
%Definici\'on del Problema (10\%):  & \underline{\hspace{2cm}} \\
%Estado del Arte (40\%):  & \underline{\hspace{2cm}} \\
%Modelo Matem\'atico (20\%): &  \underline{\hspace{2cm}}\\
%Conclusiones (15\%): &  \underline{\hspace{2cm}}\\
%Bibliograf\'ia (5\%): & \underline{\hspace{2cm}}\\
% &  \\
%\textbf{Nota Final (100\%)}:   & \underline{\hspace{2cm}}
%\end{tabular}
%---------------------------------------------------------------------------%
%\vspace{2cm}


\begin{abstract}
Abstract
\end{abstract}

\section{Modelo Matem�tico}
\subsection{Linealidad del nivel}
En una primera versi�n, la linealidad se midi� con la entrop�a de las alturas de las plataformas. Sin embargo, la entrop�a solo cuenta cu�ntas alturas distintas se usan y no c�mo est�n ordenadas, por lo que un nivel que baja de forma constante en diagonal aparece como poco lineal aunque siga una recta. Por esto se cambi� a una medida basada en la de Launchpad~\cite{Smith2011Launchpad}: se ajusta una recta a las alturas de las plataformas y se mide qu� tan lejos quedan de ella.

Sea $C$ el conjunto de columnas que tienen una plataforma y $n = |C|$ su cantidad:
\begin{equation}
C = \Big\{\, i \in I \;:\; \sum_{j \in J} x_{ijp} \geq 1 \,\Big\}
\end{equation}

Para cada columna $i \in C$, la altura de su plataforma es:
\begin{equation}
h_i = \sum_{j \in J} j \cdot x_{ijp} \qquad \forall\, i \in C
\end{equation}
Esta expresi�n supone una sola plataforma por columna, lo que se cumple en las soluciones generadas.

Luego se ajusta la recta $\hat{h}(i) = a + b\,i$ por m�nimos cuadrados, donde $\bar{i}$ y $\bar{h}$ son los promedios de las columnas y de las alturas:
\begin{equation}
b = \frac{\sum_{i \in C} (i - \bar{i})(h_i - \bar{h})}{\sum_{i \in C} (i - \bar{i})^2},
\qquad
a = \bar{h} - b\,\bar{i}
\end{equation}

La distancia promedio de las plataformas a la recta es:
\begin{equation}
d = \frac{1}{n} \sum_{i \in C} \big|\, h_i - (a + b\,i) \,\big|
\end{equation}

Finalmente, la linealidad se normaliza por la mitad del alto del mapa, que es la mayor distancia promedio posible, para que quede entre 0 y 1:
\begin{equation}
L = \max\left(0,\; 1 - \frac{d}{(H-1)/2}\right)
\end{equation}
Un valor de $L$ cercano a 1 indica que las plataformas siguen una recta (alta linealidad), mientras que un valor menor indica que el nivel sube y baja (baja linealidad).

\subsection{Funci�n objetivo}
El dise�dor indica la cantidad deseada $O_k$ de cada tipo de objeto y la linealidad deseada $L^{*}$. La cantidad colocada de cada objeto es $o_k = \sum_{i \in I}\sum_{j \in J} x_{ijk}$, y el error en las cantidades se normaliza por el total pedido:
\begin{equation}
D = \frac{\sum_{k \in K} |\, o_k - O_k \,|}{\sum_{k \in K} O_k}
\end{equation}

La funci�n objetivo combina ambos t�rminos:
\begin{equation}
\min\; Z = \alpha \cdot D + \beta \cdot |\, L - L^{*} \,|,
\qquad \alpha + \beta = 1,\quad \alpha, \beta \geq 0
\end{equation}

Como $D$ y $|L - L^{*}|$ quedan entre 0 y 1, los pesos $\alpha$ y $\beta$ indican directamente la importancia de cada parte: por ejemplo, con $\alpha = 0{,}7$ y $\beta = 0{,}3$ las cantidades pesan un 70\,\% y la linealidad un 30\,\%. Como $L$ depende de una recta ajustada y de valores absolutos, la funci�n objetivo no es lineal, por lo que el problema se resuelve con m�todos heur�sticos (Greedy y Hill Climbing) en vez de un solver exacto.

\bibliographystyle{plain}
\bibliography{Referencias}

\end{document} 